%% AR-601M arm -- kinematic scheme, drawn from the factor chain (Task 1):
%%   0T6 = Rz(t1) Tz(l1) Rz(10deg*b) Ry(t2) Ty(-l2*b) Rx(t3) Tz(-l3)
%%         Rz(t4) Ry(t5) Rx(t6) Tx(l4)
%% An original schematic, not a redraw of the source photograph. \input
%% this file from both the deck (F03) and the handout; it only defines
%% \arsixonescheme, it draws nothing on its own.
%%
%% Uses tikz libraries calc, positioning, arrows.meta -- all already
%% loaded by both documents' preambles.
\newcommand{\arsixonescheme}{%
\begin{tikzpicture}[>=Latex, line width=0.7pt, scale=0.8,
    axis/.style={<->, thick},
    offset/.style={->, thick, blue!55!black},
    frame/.style={->, thick, black!70}]
  \coordinate (O) at (0,0);
  \coordinate (A) at (0,4.2);
  \coordinate (B) at (2.6,4.2);
  \coordinate (C) at (2.6,1.6);
  \coordinate (T) at (4.6,0.4);

  % base frame, at the fixed shoulder base
  \draw[frame] (O) -- ++(0.5,0) node[right] {\scriptsize $x_0$};
  \draw[frame] (O) -- ++(0,0.5) node[above] {\scriptsize $z_0$};
  \node[below=2pt] at (O) {\scriptsize base};
  \draw[axis] ($(O)+(1.1,-0.4)$) -- ++(0,0.8) node[right] {\scriptsize $t_1\,(z)$};

  % l1: fixed offset up to the shoulder -- label to the left, clear of t1
  \draw[offset] (O) -- (A) node[midway, left=2pt] {\scriptsize $l_1$};

  % fixed 10 degree shoulder tilt about z (not a joint)
  \draw[thin] (A) ++(160:0.45) arc (160:70:0.45);
  \node[above, align=center] at ($(A)+(0,0.7)$)
    {\scriptsize $10^\circ$ tilt\\[-2pt]\scriptsize(fixed)};

  % t2 about y -- label above the arrow
  \draw[axis] ($(A)+(0.9,0)$) -- ++(0.7,0) node[above] {\scriptsize $t_2\,(y)$};
  % l2 along -y*b -- label below the A-B segment
  \draw[offset] (A) -- (B) node[midway, below=2pt] {\scriptsize $l_2\,(-y{\cdot}b)$};

  % t3 about x -- axis out of the page, drawn as a target glyph
  \draw (B) circle (2.2pt); \fill (B) circle (0.6pt);
  \node[above right] at (B) {\scriptsize $t_3\,(x)$};

  % l3: down to the wrist centre -- label to the right of the segment
  \draw[offset] (B) -- (C) node[midway, right=2pt] {\scriptsize $l_3$};

  % wrist centre: t4, t5, t6 share this one point (no offset between them)
  \fill (C) circle (1.6pt);
  \node[above left=1pt] at (C) {\scriptsize wrist centre};

  % t4 (z), t5 (y), t6 (x) fanned well clear of C, of each other, of l3
  % above and l4 below-right, and of the t1/base cluster near the origin
  \draw[axis] ($(C)+(-1.0,0.85)$) -- ++(0,0.7) node[left] {\scriptsize $t_4\,(z)$};
  \draw[axis] ($(C)+(-1.55,-0.1)$) -- ++(0.7,0) node[below] {\scriptsize $t_5\,(y)$};
  \draw ($(C)+(-0.25,-1.3)$) circle (2.2pt); \fill ($(C)+(-0.25,-1.3)$) circle (0.6pt);
  \node[below] at ($(C)+(-0.25,-1.3)$) {\scriptsize $t_6\,(x)$};

  % l4: offset to the tool frame
  \draw[offset, dashed] (C) -- (T) node[midway, below=3pt] {\scriptsize $l_4\,(x)$};
  \draw[frame] (T) -- ++(0.5,0.2) node[right] {\scriptsize $x_6$};
  \draw[frame] (T) -- ++(-0.1,0.5) node[above] {\scriptsize $z_6$};
  \node[below] at (T) {\scriptsize tool};

  % arm mirror, as a note rather than a second drawing
  \node[align=left, font=\scriptsize, text width=3.4cm] at (7.4,2.6)
    {$b=+1$ shown (right arm).\\ $b=-1$ mirrors $l_2$ and the shoulder
     tilt about $z$ (left arm).};
\end{tikzpicture}%
}
